\subsection{Existing Systems: Detailed Analysis}
\label{subsec:paperanalysis}

This section presents the state of the art in cloud-native healthcare security. Each paper is examined through three lenses: the proposed architecture, the operational mechanism, and the residual gaps that remain after the paper's contribution. Papers are grouped thematically to make the survey navigable.

\subsubsection{Front-End and Back-End Security Architectures}

{Paper 1 --- R.~L.~et al., \textit{Security in Healthcare} (ICCES 2025).} This paper proposes a dual-layered cyber security architecture that separates security concerns into front-end (portals, mobile applications, dashboards, tele-medicine interfaces) and back-end (servers, databases, cloud storage, APIs, processing engines) components. The front-end defences apply secure user authentication, encrypted data transmission, input validation, and countermeasures against interface-level vulnerabilities. The back-end applies data-at-rest encryption, access control, API security design, intrusion detection, and continuous monitoring to identify anomalies. The framework aligns with HIPAA, GDPR, and ISO 27001 and incorporates employee training to minimise insider threats. Despite its comprehensiveness as a conceptual model, the paper does not present a concrete implementation: the dual-layer separation is not coupled to any dynamic posture assessment, no offline survivability is provided, no resource-efficient edge path is proposed, no automated remediation is implemented, and no P99 tail-latency distributions are reported. The paper itself acknowledges that ``no single solution is sufficient'' but does not propose the unified architecture required to resolve that limitation.

\subsubsection{Runtime Detection Architectures}

\textbf{Paper 2 --- L.~Ali et al., \textit{Intelligent Security Framework for Real-Time Threat Detection} (ICCR 2025).} The framework combines Deep Reinforcement Learning (DRL) with Federated Anomaly Detection (FAD). FAD uses LSTM autoencoders and One-Class Support Vector Machines (OC-SVM) to detect zero-day attacks and Advanced Persistent Threats locally; DRL, implemented with both DQN and PPO, learns optimal response actions through trial-and-error interaction with the environment. Model updates aggregate in the cloud without transferring raw patient data, preserving privacy. However, the DRL pipeline with LSTM autoencoders is computationally intensive: the paper reports only mean performance metrics and does not test on edge-class hardware. There is no static posture management, no dynamic posture assessment, no offline survivability (the architecture assumes continuous cloud connectivity for model aggregation), and no automated remediation. Data residency and sovereignty are not addressed despite being relevant to the healthcare deployments the paper targets.

{Paper 4 --- P.~J.~Gunjale et al., \textit{Intelligent Cloud Security Model for Healthcare} (ICCTWC 2026).} The paper proposes an ML-based threat monitoring framework with continuous compliance assurance. The system analyses cloud system logs, network traffic, and user behaviour in real time to detect anomalous activities while enforcing HIPAA and GDPR through automated compliance monitoring. The paper reports 94\% detection accuracy and reduced false positives against a conventional rule-based security system. The evaluation, however, is conducted entirely in a simulated cloud environment, not on physical constrained hardware. There is no dynamic posture management---compliance is checked periodically rather than continuously---and no static posture auditing is integrated with the runtime detection. The paper reports accuracy but omits P99/P95 tail-latency distributions, and no automated remediation is implemented.

{Paper 7 --- A.~R.~Abusalih Rahumath Ali et al., \textit{Enhancing Cloud Security with AI} (ICSCNA 2024).} The paper investigates the application of machine learning for cloud security with anomaly detection, behaviour analysis, and automated incident response. The framework employs ML algorithms and advanced analytics to proactively discover and respond to threats in real time, and presents case studies of successful AI-enhanced security deployments. The evaluation uses simulated environments only, not physical constrained hardware. No static or dynamic posture management is implemented, no offline survivability is provided, and no P99/P95 tail-latency distributions are reported. The automated incident response is discussed conceptually but not implemented.

{Paper 15 --- B.~M.~and K.~M.V., \textit{Real-Time Threat Detection in Cloud-Based Telemedicine} (ICCDS 2025).} The paper applies Deep Reinforcement Learning to cloud telemedicine, with an intelligent agent that learns optimal defence policies through environment interaction. The system adapts to ransomware, session hijacking, and inappropriate access, and the authors report 98.5\% detection accuracy, 35\% improvement in response time, and 2.1\% false positive rate. The DRL architecture is very expensive computationally and is not tested on edge hardware. No static or dynamic posture management is implemented; the evaluation uses simulated healthcare attack logs rather than physical deployment; no P99/P95 tail-latency distributions are reported; automated remediation exists inside the DRL policy but is not integrated with posture management.

{Paper 17 --- P.~Joshi and S.~Singh, \textit{TrishulFL} (Cyber-AI 2025).} TrishulFL is a comprehensive security framework for federated learning in healthcare. It integrates real-time threat detection, attack mapping, automated incident response, and self-healing. The mechanism combines real-time anomaly detection, spatio-temporal threat visualisation, reinforcement learning agents, and generative AI explainability. The framework is evaluated only in controlled simulations, requires high computational resources for RL and generative AI components, and assumes continuous connectivity for federated learning coordination. No static or dynamic posture management is implemented, no offline survivability is provided, and no P99/P95 tail-latency distributions are reported. Remediation exists within the framework but is not integrated with posture management.

{Paper 23 --- H.~H.~Bin Suhaimi and H.~T.~Zubair, \textit{Data Leakage Detection in Cloud Computing Environment} (CyberComp 2024).} The paper evaluates data leakage detection methods with a focus on insider threats, using Random Forest and k-Nearest Neighbours for anomaly detection and behavioural analytics. The evaluation uses real-world data but not physical constrained hardware. No posture management is implemented, no offline survivability is provided, no P99/P95 tail-latency distributions are reported, and no automated remediation is implemented. Notably, the framework decouples detection from encryption and access control, so the detection logic cannot respond to enforced policies.

\subsubsection{Blockchain and Integrity-Centric Architectures}

{Paper 3 --- G.~Thakre et al., \textit{AI-Enhanced Security in Blockchain and Cloud-Based EHR} (ICMLAS 2025).} This paper reviews the combination of Artificial Intelligence, Blockchain, and cloud computing for Electronic Healthcare Records. Blockchain provides immutable data integrity and privacy through a distributed ledger; cloud platforms provide scalable infrastructure; AI provides threat identification and anomaly detection. The paper discusses quantum-resistant algorithms, federated learning, and AI-powered smart contracts. It is a review, not an implementation: no concrete architecture is presented, blockchain consensus overhead makes the design unsuitable for constrained edge devices, no runtime anomaly detection is implemented, no dynamic posture management is provided, and no performance metrics are reported.

{Paper 5 --- S.~D.~et al., \textit{Blockchain-Enhanced Hybrid Cloud Healthcare Data Management} (ICERCS 2025).} The paper proposes a blockchain-backed hybrid cloud combining the Google Cloud Healthcare API, Hyperledger Fabric, and AWS Health Lake. It achieves 74\% latency reduction, 34\% security improvement, and 228\% throughput increase, with greater than 94\% accuracy on image analysis and disease prediction. Smart contracts enforce compliance and blockchain audit trails are immutable. However, the blockchain consensus imposes high computational and bandwidth overhead; the paper does not validate on physical constrained hardware, implements no runtime anomaly detection, provides no dynamic posture management, assumes continuous cloud connectivity, reports no P99/P95 distributions, and does not implement automated remediation.

\subsubsection{Multi-Layer Zero-Trust and Hybrid Frameworks}

{Paper 6 --- K.~Yadav et al., \textit{Security in Healthcare: Protecting Sensitive Patient Data} (IBI 2026).} The paper presents a multi-layered framework combining Zero Trust Network Access (ZTNA), encrypted IoMT gateways, blockchain-supported audit logs, federated identity management, and a compliance engine. Zero-trust authentication applies to every request regardless of source; the compliance engine enforces HIPAA and GDPR through automated policy checks. The paper is conceptual and provides no empirical validation. It assumes continuous connectivity for ZTNA and blockchain, provides no offline survivability, and only partially addresses dynamic posture management (through ZTNA, but without continuous risk scoring or adaptive policy updates). The framework requires significant computational resources for encryption, blockchain operations, and identity management, and no P99/P95 tail-latency distributions are reported.

{Paper 13 --- R.~G.~S.~and S.~H.~J., \textit{Ensuring Secure Healthcare Data Storage with ANFIS and CP-ABE} (ICIMIA 2025).} The paper integrates Ciphertext-Policy Attribute-Based Encryption (CP-ABE) with an Adaptive Neuro-Fuzzy Inference System (ANFIS). CP-ABE provides fine-grained access control via access policies embedded in ciphertext; ANFIS tracks access behaviour and predicts threatening events. CP-ABE is roughly 100$\times$ slower than AES-256 on ARM processors, which is a fundamental barrier to edge deployment. The framework focuses exclusively on encryption and access control, not on runtime anomaly detection; it provides no offline survivability; it uses simulations only; it reports no P99/P95 tail-latency distributions; and automated remediation is not implemented.

{Paper 18 --- R.~Kande, \textit{Z-HealthcareAI} (ICNWC 2026).} Z-HealthcareAI enhances zero-trust security for healthcare clouds with behavioural identity analytics, federated anomaly detection, encrypted data collection, and policy-driven compliance automation following NIST SP 800-207. Behavioural biometrics, device validation signals, and network anomalies feed identity risk scores. The framework was tested using MIMIC-III, CIC-IDS 2018, and MedBIoT datasets with OMNeT++ simulation. It requires high computational resources for AI-based identity analytics and federated anomaly detection; it assumes continuous connectivity for federated learning and compliance enforcement; it does not implement static or dynamic posture management; it provides no offline survivability; it reports no P99/P95 tail-latency distributions; and no automated remediation is implemented.

\subsubsection{Federated and Privacy-Preserving Detection}

{Paper 8 --- J.~Simon and N.~Kapileswar, \textit{Federated Deep Learning-Driven Cloud-IoT Framework} (ICSSAS 2025).} The paper proposes a Federated Deep Learning (FDL) Cloud-IoT framework for real-time healthcare monitoring and privacy-preserving anomaly detection. IoT-enabled wearable devices train local deep learning models, and only model updates are aggregated in the cloud. A lightweight anomaly detection algorithm detects arrhythmias and sudden falls. The framework requires high compute at local devices; it assumes continuous connectivity for aggregation; it provides no static or dynamic posture management; it offers no offline survivability; it uses benchmark datasets only; it reports no P99/P95 distributions; and it does not implement automated remediation.

{Paper 10 --- N.~Kapileswar and J.~Simon, \textit{PrivAudit} (ICICI 2025).} PrivAudit uses Federated Transformers with Differential Privacy to detect cloud anomalies. Transformers capture temporal and spatial patterns in network traffic; federated learning keeps data local; differential privacy formalises gradient protection. The Transformer architecture requires very high memory and CPU resources; the framework assumes continuous connectivity for model aggregation; it provides no static or dynamic posture management; it offers no offline survivability; it uses standard cloud security datasets rather than constrained hardware; it reports no P99/P95 distributions; and it does not implement automated remediation.

{Paper 12 --- K.~Joshi et al., \textit{Secure and Privacy-Preserving Healthcare Data Management} (MISSF 2026).} The paper proposes an AI-powered intrusion detection model that integrates federated deep learning with differential privacy for distributed clouds. It trains on encrypted datasets, achieves 98.3\% detection accuracy, and reduces false alarms by 3.7\%. The framework requires high compute, assumes continuous connectivity, and provides no static or dynamic posture management. It offers no offline survivability, uses benchmark data only, reports no P99/P95 distributions, and does not implement automated remediation.

\subsubsection{Edge and IoMT Architectures}

{Paper 11 --- P.~Arunachalam et al., \textit{Applications of IoMT} (ICBDS 2025).} The paper explores IoMT applications for intelligent health environments, combining wearable sensors, medical devices, cloud computing, and AI. Sensors stream to cloud-hosted platforms; AI identifies anomalies and predicts health decline; telemedicine integration follows HL7 and FHIR; encryption protocols and edge computing are discussed. The architecture assumes continuous cloud connectivity, offers no offline survivability, implements no posture management, and reports only clinical anomalies rather than security anomalies. The evaluation uses case studies and simulations, reports no P99/P95 distributions, and does not implement automated remediation.

{Paper 14 --- M.~R.~and G.~Keerthana, \textit{Developing IoT Healthcare Networks Based on Big Data Analytics} (IDCIoT 2026).} The paper proposes an intelligent healthcare model using AWS IoT Core, Kinesis Data Analytics, and Health Lake with Digital Twins. MQTT carries medical data; FHIR carries medical documents; an Advanced Integrated Data Security (AIDS) framework layers encryption, access control, anomaly detection, and intrusion prevention. The architecture assumes continuous cloud connectivity, provides no offline survivability, does not implement dynamic posture management, and is conceptual rather than empirically validated. No P99/P95 distributions are reported and no automated remediation is implemented.

\textbf{Paper 19 --- R.~Katiyar et al., \textit{Application of AI-Driven Techniques for Smart Healthcare Monitoring} (ICICIT 2026).} The paper proposes an AI-based framework synergising IoT sensor analytics, machine learning, and cloud-edge hybrid computing. It addresses data heterogeneity, latency, privacy, and interoperability through dynamic task allocation. The framework requires high compute, implements no posture management, focuses on clinical monitoring rather than security, provides no offline survivability, uses simulations only, reports no P99/P95 distributions, and does not implement automated remediation.

{Paper 20 --- R.~U.~Islam et al., \textit{NDN-Aided Edge-Cloud Architecture} (FiCloud 2025).} The paper exploits Named Data Networking (NDN) with a hybrid edge-cloud paradigm for continuous smart healthcare. NDN provides efficient data-named communication; the edge-cloud split balances local processing with cloud services. This is the only paper in the corpus with offline survivability as a first-class property. However, it does not provide a unified security framework, does not implement posture management, does not implement runtime security detection, uses OMNeT++ simulation only, reports no P99/P95 distributions, and does not implement automated remediation.

{Paper 22 --- Affrose et al., \textit{Health Cloud-Based IoT Virtual Health Monitoring} (ICCES 2024).} The paper presents an ESP8266-based health monitoring system using ECG (AD8232), DS18B20 temperature, MAX30100/MAX30102 pulse oximeter, BMP280 respiration rate, and blood pressure sensors. Sensors stream to the Arduino Cloud for remote access. This is the only paper in the corpus with hardware validation. It assumes continuous cloud connectivity, provides no offline survivability, implements no posture management, implements no runtime security detection, reports no P99/P95 distributions, and does not implement automated remediation.

{Paper 21 --- R.~J.~Shimpi and V.~Tiwari, \textit{PRE-MWOA} (ICDSAAI 2025).} The paper introduces Proxy Re-Encryption with a Modified Whale Optimization Algorithm for Health IoT in fog computing. PR protects data during task distribution; MWOA optimises anomaly detection; late-stage decryption protects confidentiality. The framework requires moderate compute; it implements no posture management; it focuses on encryption; it provides no offline survivability; it uses simulations only; it reports no P99/P95 distributions; and it does not implement automated remediation.

\subsubsection{Encryption and Access-Centric Systems}

{Paper 25 --- R.~T.~Dinesh Kumar et al., \textit{Secure Medical Record Storage using Deep Encryption} (ICEARS 2026).} The paper proposes Distributed Customizable Unique Node Access (DCUNA/CUNA), which segments medical images and encrypts each segment across distributed nodes with unique keys and access policies. The framework is implemented in Java for platform independence. The approach requires moderate compute; it implements no posture management; it focuses entirely on encryption; it provides no offline survivability; it uses simulations only; it reports no P99/P95 distributions; and it does not implement automated remediation.

\subsubsection{Reviews and Surveys}

{Paper 9 --- T.~K.~Kota, \textit{Multi-Cloud Data Engineering for GenAI} (ASSIC 2025).} The paper proposes a multi-cloud data engineering framework for GenAI-driven medical insurance analytics. It combines multi-cloud storage, serverless computing, and federated learning; foundation models improve ingestion, transformation, and predictive analytics. The framework requires significant compute; it implements no posture management; it is not security-focused; it provides no offline survivability; it uses simulated environments; it reports no P99/P95 distributions; and it does not implement automated remediation.

{Paper 16 --- G.~Narsimlu et al., \textit{Assessment of AI in Hybrid Cloud Security} (ICECA 2025).} The paper is a review study classifying deep learning algorithms for cloud security and categorising threats. It investigates AI-driven adaptive risk assessment. Being a review, it presents no implementation, no posture management, no offline capability, no resource-constrained consideration, no P99/P95 distributions, and no automated remediation.

{Paper 24 --- G.~Nithyavani and G.~Naga Raja, \textit{A Comprehensive Survey on Security and Privacy in IoMT} (IEEE Access 2025).} The paper is a survey reviewing IoMT security with an attack taxonomy, classification of protocols, and future research directions. Being a survey, it presents no implementation, no posture management, no offline capability, no resource-constrained consideration, no P99/P95 distributions, and no automated remediation.

